\section{Confidential Inference}

\label{sec:fhe}
% ============================================================================

Hanzo implements confidential inference using the CKKS (Cheon--Kim--Kim--Song)
FHE scheme~\cite{ckks2017}, which supports approximate fixed-point arithmetic
on encrypted vectors---the operation profile that neural network inference
requires.

\subsection{CKKS Overview}

CKKS operates on the ring $R_Q = \mathbb{Z}_Q[X]/(X^N + 1)$ where $N$ is a
power of 2 (typically $N = 2^{16}$) and $Q$ is a product of coprime moduli.
A plaintext vector $\bm{z} \in \mathbb{C}^{N/2}$ is encoded as a ring
element $m \in R_Q$ via the canonical embedding, then encrypted as:

\begin{equation}
\text{ct} = (\bm{a} \cdot s + e + m, \; \bm{a})
\end{equation}

where $s$ is the secret key, $\bm{a}$ is uniform random, and $e$ is a small
error. Homomorphic addition and multiplication on ciphertexts correspond to
addition and multiplication on plaintexts, with the error growing after each
operation.

\subsection{Inference Pipeline}

The confidential inference pipeline operates in four phases:

\begin{enumerate}[leftmargin=1.4em]
    \item \textbf{Client-side encryption}: The client encodes the input (token
    embeddings for text, pixel values for images) into CKKS plaintext vectors
    and encrypts them with their public key. The ciphertexts are transmitted
    to the Hanzo Cloud.
    \item \textbf{Homomorphic evaluation}: The cloud evaluates the model on
    the ciphertexts. Linear layers (matrix multiplication, convolution) map
    directly to CKKS operations. Non-linear activations (GELU, softmax) are
    approximated by low-degree polynomials fitted via Remez
    algorithm~\cite{hanzocandle, hanzofhe}.
    \item \textbf{Encrypted output}: The cloud returns the encrypted result.
    It cannot decrypt it---only the client (or threshold decryption
    participants) hold the secret key.
    \item \textbf{Client-side decryption}: The client decrypts the result
    locally. For threshold configurations, $k$-of-$n$ key holders cooperate
    to produce the decryption.
\end{enumerate}

\subsection{Polynomial Activation Approximation}

The GELU activation $\text{GELU}(x) = x \cdot \Phi(x)$ where $\Phi$ is the
Gaussian CDF is approximated by a degree-7 polynomial on $[-8, 8]$:

\begin{equation}
\widetilde{\text{GELU}}(x) = \sum_{k=0}^{7} a_k x^k
\end{equation}

with coefficients fitted to minimize $L_\infty$ error. The maximum
approximation error is $|\text{GELU}(x) - \widetilde{\text{GELU}}(x)| <
2 \times 10^{-4}$ on the domain, which is within the noise floor of CKKS
arithmetic. Softmax is decomposed into exponentiation (degree-6 polynomial
approximation) and reciprocal (Goldschmidt iteration on ciphertexts).

\subsection{Threshold Decryption}

For multi-party scenarios, the secret key $s$ is split into $n$ shares
$s_1, \ldots, s_n$ via Shamir's secret sharing. Decryption requires $k$
shares:

\begin{equation}
m = \text{ct}_0 + \text{ct}_1 \cdot s = \text{ct}_0 +
\text{ct}_1 \cdot \sum_{i \in S} \lambda_i s_i
\end{equation}

where $S \subset [n]$ with $|S| = k$ and $\lambda_i$ are Lagrange
interpolation coefficients. Each participant computes a partial decryption
$d_i = \text{ct}_1 \cdot \lambda_i s_i + e_i$ (with a smudging noise $e_i$
to prevent key leakage), and the combiner sums the partial decryptions. No
single party learns the plaintext unless they hold $k$ shares.

% ============================================================================
